%
%============= CHAPTER BREAK =================%
%
\chapter{Artemidorus of Daldis (ca 180--200 CE)}%
\vspace*{3mm}
\justifying
%
\noindent{}Largely known for his \textit{Oneirocritica}, %
Artemidorus of Daldis likely wrote during the late Second Sophistic. %
He dedicated numerous books of his extensive work on dream interpretation %
to his intellectual contemporary, Maximus of Tyre, %
and makes mention of several other persons known to have lived between the 170s CE and %
the early Severan period.\endnote{Bowersock 2004:54--56.%
} %
Enormously influential in modern psychology and dream-analysis, Artemidorus is hardly ever %
considered as a potential contemporary to New Testament author(s)/editor(s). %
But in our view, this is prejudicial, and could get in the way of observing %
historically and literarily secure connections between his works %
and late editorial stages of New Testament documents. \\[1.5mm] %%

\noindent{}\textbf{ἀνένδεκτος} / \textit{anendektos} / “inconceivable” % 
appears in two results tied to BCE authors, %
both found within  modern compilations of ancient fragments: 
Chrysippus, \textit{Frag}. 55 and Aristophanes Gram., \textit{Frag}. 13. %
The second is clearly part of the secondary wording of Eustathius of Thessalonica (12th cent. CE), %
who subsequently cites to Chrysippus only in regard to diminutive forms of address for beloved relatives and the like. %
The other fragment requires a more careful examination. %
It is annotated as “Diocles Magnes apud Diog. Laërt. VII, 50”, %
i.e., Diocles the Magnesian within Diogenes Laertius 7.50. %
Diogenes Laertius wrote in the early 3rd century CE, %
but Diocles of Magnesia is understood to have lived in the 2nd or 1st century BCE. %
In the context of the fragment, the lemma is attributed to Diocles of Magnesia, %
but since we lack any works authored by Diocles to test that attribution, %
the attestation remains dubious. %
 %
\clearpage
 \begin{quote}\textls[-15]{διαφέρει δὲ φαντασία καὶ φάντασμα. %
 φάντασμα μὲν γάρ ἐστι δόκησις διανοίας, οἵα γίνεται κατὰ τοὺς ὕπνους, %
 φαντασία δέ ἐστι τύπωσις ἐν ψυχῇ, τουτέστιν ἀλλοίωσις, %
 ὡς ὁ Χρύσιππος ἐν τῇ βʹ περὶ ψυχῆς ὑφίσ\-ταται. %
 οὐδὲ γὰρ δεκτέον τὴν τύπωσιν οἱονεὶ τύπον σφραγιστῆρος, %
 ἐπεὶ \textbf{ἀνένδεκτόν} ἐστι πολλοὺς τύπους κατὰ τὸ αὐτὸ περὶ τὸ αὐτὸ γίνεσθαι.} \\[1mm] %
 Now fantasy and phantasm are distinct. %
 For phantasm is a fancy of mind, such as occurs according to dreams, %
 but fantasy is an impression in soul, that is, an alteration, %
 as Chrysippus substantiates in book two \textit{On the Soul}. %
 Neither should one understand the impression as akin to the mold of a signet, %
 since it is \textbf{inconceivable} that many molds based on the same thing %
 be applied around the same thing. %
 \end{quote} % 
 \vspace*{2mm}

 Thus the first secure non-Christian attestation of this lemma, %
 both in terms of chronology and authorship, %
 is Artemidorusʼ \textit{Oneirocritica}. %
 In context, the author discusses how to interpret numbers, %
 as well as the numerical value of words (based on the number of letters in the word). %
 All are calculated or adjusted to correlate with typical human life-spans and %
 reasonable periods of time remaining in a personʼs life. %

\vspace*{2mm} 
 \begin{quote}ὁ δὲ παραπλήσιος καὶ ἐπὶ τῶν ἄλλων ὁμοίων λόγος τηρείσθω,
 ὅταν ᾖ ὁ λεγόμενος ἀριθμὸς τοῦ μὲν ἔμπροσθεν χρόνου ἥττων, %
 πρὸς δὲ τὸν μέλλοντα \textbf{ἀνένδεκτος}. (\textit{Oneir}. 2.70) \\[1mm]
 One must apply the equivalent logic also to other similar ones, %
 whenever a spoken number be less than the preceding time, %
 but would be \textbf{inconceivable} for the time to come. \end{quote} % 
 \vspace*{2mm}
 %
The solitary instance in canonical NT texts is found in Luke 17:1, %
where the term is missing from close parallels with canonical Matthew, %
canoncial Mark, and Marcion's \textit{Evangelion}. %

\vspace*{2mm} 
 \begin{quote}\textbf{Ἀνένδεκτόν} ἐστιν τοῦ τὰ σκάνδαλα μὴ ἐλθεῖν, % 
  πλὴν οὐαὶ δι’ οὗ ἔρχεται· / \textbf{Inconceivable} it is that the scandals not come, %
 however, woe [to him] through whom it comes. %
\end{quote} % 
 \vspace*{2mm}

The woe statement not only shares a new/rare lemma with Artemidorus, %
but also precedes a clear evocation of heinous death (Luke 17:2). %
In its context, it thus aligns with Artemidorusʼ thematic concern about the proper interpretation of dreams %
in relation to human mortality. \\[1.5mm] %

\noindent{}\textbf{βραδυπλοέω} / \textit{braduploeō} / “to sail slowly”, %
“to drift”, or “to be in the doldrums” has only 11 TLG hits and appears first in non-Christian writings %
in the \textit{Oneirocritica}. %
As with the rest of the work, %
the focus here is on the interpretation and meaning of dreams. %

\vspace*{3mm} 
 \begin{quote}Πάντα τὰ περιεκτικὰ ἡμῶν αὐτῶν τὸν αὐτὸν ἔχει %
 λόγον, ἱμάτιον οἰκία τεῖχος ναῦς καὶ τὰ ὅμοια. οἷον ξύλινον %
 ἱμάτιον ἔδοξέ τις περιβεβλῆσθαι. %
 ἔτυχε πλέων καὶ \textbf{ἐβραδυπλόησεν}· % 
 ἦν γὰρ αὐτῷ ξύλινον ἱμάτιον ἡ ναῦς. (\textit{Oneir}. 4.30) \\[1mm]
\textls[-10]{All things that beset our own selves have the same rationale: %
 garment, house, wall, ship, and the like. Someone thought %
 he was enclothed in some sort of wooden garment. It chanced %
 that while sailing \textbf{he sailed slowly}. %
 For the ship was a wooden garment to him.} %
 \end{quote} % 
 \vspace*{3mm}

 If the influence of Tolkienʼs mythical tropes on our thinking isnʼt yet apparent, %
 let us rephrase a line from the \textit{Fellowship of the Ring}. %
 For Artemidorus, the dream image of %
 a ship as a wooden garment portended bodily ill at sea and conveyed a sense of Moria-like entrapment: %
 “This isnʼt a ship. Itʼs a casket”. %
Artemidorusʼ mention of “sailing slowly” in connection with an %
individualʼs foreboding dream about a dangerous sea-voyage is fascinating %
to consider on its own, but even more fascinating to compare with the solitary %
use of the same unusual term in the canonical NT. %
It appears near the start of Paulʼs extended sea-voyage to Italy (Acts 27), %
accentuating a slow stretch of sailing near and %
leading to Knidos (Acts 27:6).

\vspace*{3mm} 
 \begin{quote}ἐν ἱκαναῖς δὲ ἡμέραις \textbf{βραδυπλοοῦντες} καὶ μόλις γενόμενοι κατὰ τὴν Κνίδον / %
 Now \textbf{sailing slowly} across quite enough days, and with difficulty coming down to Knidos\end{quote} % 
 \vspace*{3mm}

\noindent{}Knidos was a port city in southwest Asia Minor famed for the ancient Battle of Knidos (394 BCE). %
It was commonly associated with Artemidorus of Knidos (not Daldis). %
This Artemidorus was celebrated by Plutarch (\textit{Lives}, Caesar 65) as loyal friend to Julius Caesar. %
The local games were held in his honor, and he was ritually deified after his death. %
One Artemidorus deserves another, it seems. %

The comparisons become more intriguing in the literary context of Acts, %
since Paul soon thereafter delivers an evil omen (Acts 27:10). %

\vspace*{3mm} 
 \begin{quote}ἄνδρες, θεωρῶ ὅτι μετὰ ὕβρεως καὶ πολλῆς ζημίας οὐ μόνον %
    τοῦ φορτίου καὶ τοῦ πλοίου ἀλλὰ καὶ τῶν ψυχῶν ἡμῶν μέλλειν ἔσεσθαι τὸν πλοῦν. / %
 Men, I perceive that the voyage is about to meet with wanton violence and much damage %
 not only to the cargo and the ship, %
 but also to our very lives.\end{quote} % 
\vspace*{3mm}

\noindent{}This omen is said to differ from the professional opinion of the shipʼs captain and pilot (27:11). 
As the story in Acts proceeds, Paul not long after has near-death experiences at sea and on the island of Malta. %

So, Acts aligns with Artemidorus in several respects, %
not only a rare and perhaps newly coined and/or popularized term, %
but also a foreboding omen of dangerous travel at sea. %
The distinctive parallels multiply on broader consideration. %
Acts uses this term just before mentioning Knidos, %
a city famed for its empire-supporting son, Artemidorus, %
the same name as the writer in whose work this lemma first appears. %
Artemidorus was a famous friend of Julius Caesar who informed %
against Caesarʼs enemies in regard to an assassination plot. %
The Paul of Acts is expressly on a journey for a personal audience with Caesar that he requested. %
Is this all mere coincidence, or is it brilliant literary allusion and strategy %
on the part of the editor of canonical Acts? %
If the latter, then this points to the sea-voyages in Acts 27 as very late-2nd or, more likely, early-3rd century creation. %
%

